\documentclass[class=report, crop=false]{standalone}
\usepackage[T1]{fontenc}
\usepackage[utf8]{inputenc}
\usepackage{amsmath}
\usepackage{amssymb}
\usepackage{bm}
\usepackage[authoryear, round]{natbib}

\begin{document}

\chapter{Introduction}

\section{Introduction to Light Scattering}

Light scattering is a ubiquitous and fundamental physical phenomenon that describes the interaction of electromagnetic waves with discrete particles or inhomogeneous media. Whether it is the blue color of the sky, the brilliance of a white cloud, or the operation of radar systems, scattering occurs whenever an incident electromagnetic wave encounters an obstacle. This interaction induces oscillating multipoles, described as microscopic electric charge separations that vary with time within the object, which subsequently radiate secondary waves known as the scattered field. Understanding how the amplitude, phase and polarization of light evolve during this interaction is essential for a wide array of scientific and engineering disciplines.

\vspace{12pt}

The detailed study of these interactions is a cornerstone for the development of advanced technologies. In remote sensing, the analysis of scattered waves allows for the detection and characterization of objects that are otherwise inaccessible. Similarly, in biomedical optics, the scattering properties of tissues are exploited to diagnose diseases at an early stage. However, the complexity of practical scatterers, which often possess irregular shapes and inhomogeneous material properties, renders analytical solutions such as the Mie series insufficient \citep{Wriedt2012}. While Mie theory provides an exact mathematical solution for scattering by simple spheres and has served as the gold standard for validation in the field, it cannot be applied to the arbitrary geometries encountered in practical engineering problems. Consequently, numerical simulation has become an indispensable tool. By discretizing the governing equations of electromagnetism and converting the continuous differential equations into a system of algebraic equations solvable by computer, numerical methods allow researchers to predict scattering behavior with high fidelity \citep{Rumpf2014}.

\vspace{12pt}

The physics of this interaction is governed entirely by the laws of classical electromagnetism. At the macroscopic level, the behavior of these electromagnetic fields is described by a set of coupled partial differential equations known as Maxwell's equations.

\subsection{Maxwell's Equations}

In the context of light scattering, the fundamental behavior of electromagnetic fields is described by Maxwell's equations. For a linear, isotropic medium characterized by permittivity $\epsilon$ (a measure of electric polarization), permeability $\mu$ (a measure of magnetic response) and conductivity $\sigma$, the general time-domain differential forms of these equations are:

\begin{align}
\nabla \times \bm{E} &= -\frac{\partial \bm{B}}{\partial t} \label{eq:faraday} \\
\nabla \times \bm{H} &= \frac{\partial \bm{D}}{\partial t} + \bm{J} \label{eq:ampere} \\
\nabla \cdot \bm{D} &= \rho \label{eq:gauss_e} \\
\nabla \cdot \bm{B} &= 0 \label{eq:gauss_m}
\end{align}

\noindent where $\bm{D} = \epsilon \bm{E}$ is the electric flux density and $\bm{B} = \mu \bm{H}$ is the magnetic flux density.

\vspace{12pt}

\noindent\textbf{Faraday's Law of Induction \eqref{eq:faraday}.}

\noindent This equation describes the fundamental principle of electromagnetic induction. It states that a time-varying magnetic flux density ($\bm{B}$) induces a spatially circulating electric field intensity ($\bm{E}$) \citep{Jin2015}. Physically, this coupling implies that a changing magnetic environment is capable of generating an electromotive force, which is the mechanism governing devices like transformers and generators \citep{Jin2015}.

\vspace{12pt}

\noindent\textbf{Ampère's Circuital Law \eqref{eq:ampere}.}

\noindent This law establishes that the curl of the magnetic field intensity ($\bm{H}$) is generated by two distinct sources which are the flow of conductive electric current density ($\bm{J}$) and the time rate of change of the electric flux density ($\bm{D}$) \citep{Jin2015}. The latter term, $\frac{\partial \bm{D}}{\partial t}$, represents the ``displacement current,'' which is critical for explaining how magnetic fields can exist in regions without physical charge flow (such as vacuum or dielectrics), thereby allowing for the propagation of electromagnetic waves \citep{Jin2015}.

\vspace{12pt}

\noindent\textbf{Gauss's Law for Electric Fields \eqref{eq:gauss_e}.}

\noindent Gauss's law for electric fields states that the divergence of the electric flux density ($\bm{D}$) is directly proportional to the volume electric charge density ($\rho$) \citep{Jin2015}. Physically, this signifies that electric charges act as the sources (positive charge) or sinks (negative charge) of the electric field where the flux lines diverge outwards from positive charge and converge inwards towards negative charge \citep{Jin2015}.

\vspace{12pt}

\noindent\textbf{Gauss's Law for Magnetic Fields \eqref{eq:gauss_m}.}

\noindent This equation asserts that the divergence of the magnetic flux density ($\bm{B}$) is always zero \citep{Jin2015}. Physically, this implies the non-existence of isolated magnetic charges (magnetic monopoles). Consequently, magnetic field lines are solenoidal, meaning they always form continuous, closed loops with no distinct beginning or end \citep{Jin2015}.

\subsection{Significance of Light Scattering Simulation}

The simulation of light scattering has become indispensable in biomedical engineering driven by the critical need for non-invasive diagnostic tools and the ethical constraints associated with direct human or animal experimentation \citep{Periyasamy2017, Krasnikov2019}. Recognized as the gold standard for modeling light transport, Monte Carlo (MC) simulations allow researchers to rigorously study photon interactions within complex, heterogeneous biological structures which ranging from the intricate layers of the human eye to whole-body small animal models where traditional analytical solutions often fail \citep{Periyasamy2017, Krasnikov2019}. The implementation of these methods has been revolutionized by the integration of realistic mesh-based geometries and parallel processing technologies like Graphics Processing Units (GPUs), which address the massive computational demands required to trace millions of photon paths for statistical accuracy \citep{Periyasamy2017, Krasnikov2019}. Consequently, these advanced simulations provide a robust platform for optimizing novel imaging systems, such as Optical Coherence Tomography (OCT) and Raman spectroscopy and are essential for quantifying light distribution to ensure safety in therapeutic applications like optogenetics without the risks of invasive procedures \citep{Krasnikov2019}.

\vspace{12pt}

Parallel to these advancements in the biomedical domain, light scattering simulation stands as a foundational pillar in the rapidly evolving field of nanophotonics, where modern systems involve intricate light-matter interaction regimes that extend well beyond the diffraction limit \citep{Gallinet2015}. As research progresses from simple plasmonic nanostructures to increasingly sophisticated multiscale devices, traditional analytical methods become insufficient, rendering numerical modeling central to the understanding, control and design of these complex systems \citep{Gallinet2015}. Such simulations are particularly indispensable for engineering exotic metamaterials including sub-wavelength focusing lenses and for optimizing photonic bandgap structures, where the precise manipulation of light propagation, localization and scattering is critical for performance \citep{teixeira2023finite}. Consequently, the continuous demand for highly efficient numerical tools is driven by the necessity to accurately model these elaborate physical phenomena, thereby bridging the gap between theoretical nanophotonic concepts and the engineering of functional optical devices for applications ranging from sensing and light harvesting to optical signal processing \citep{Gallinet2015, teixeira2023finite}.

\vspace{12pt}

The development of advanced biosensing platforms is driven by the critical clinical need for non-invasive, continuous glucose monitoring to manage the rising prevalence of diabetes. \citet{Hu2023} highlight that conventional finger-prick methods are painful, intermittent, and prone to patient non-compliance, while enzymatic sensors often suffer from environmental instability and poor longevity. To address these limitations, the Poly(ethylene glycol) diacrylate (PEGDA) based biosensor offers a robust alternative, operating on a chemo-mechanical transduction principle where phenylboronic acid (PBA) motifs induce a volumetric phase transition upon glucose binding \citep{Hu2023, Cai2021}. However, the mechanism of this sensor relies on a continuously deforming domain rather than a static geometry, creating a complex coupled physics scenario where mechanical deformation and optical properties are intrinsically linked. This specific application serves as a compelling archetype to demonstrate the indispensability of light scattering simulation, which transcends simple visualization to provide a fundamental method for interpreting the light-matter interactions that drive next-generation stimuli-responsive diagnostics.

\vspace{12pt}

The necessity of this simulation is further elucidated by the requirement to resolve the time-dependent anisotropic deformation inherent to laser-fabricated polymeric architectures. In the context of the PEGDA biosensor, \citet{Ennis2023} observe that microstructures fabricated via Two-Photon Polymerization (2PP) do not swell uniformly but exhibit complex warping and buckling constrained by substrate adhesion. This physical reality renders analytical models like Mie theory insufficient, as they cannot represent such distorted, low-contrast boundaries \citep{Ennis2023}. Light scattering simulation becomes the essential interpretative engine to map these realistic deformation fields into the optical domain. It provides the unique capability to decouple competing variables which is isolating the effects of geometric expansion from refractive index reduction—thereby yielding mechanistic insights that are otherwise inaccessible through experimental characterization alone.

\vspace{12pt}

Finally, this biosensor example highlights the transformative impact of simulation in enabling predictive engineering and rigorous performance benchmarking. Adopting the framework established for photonic crystal sensors, the simulation allows researchers to generate computational calibration curves that link mechanical changes to precise optical shifts, predicting critical metrics such as Limit of Detection (LOD) and Sensitivity. This facilitates computational design optimization, where nanostructural parameters are optimized in silico to maximize electromagnetic field confinement prior to physical fabrication, a strategy validated in the development of SERS substrates. By successfully replicating the field enhancements theoretically expected in these complex geometries, the simulation establishes itself as a rigorous, physically accurate methodology essential for designing and validating elastomeric devices with clinical precision \citep{Kumar2025}.

\section{Problem Statement}

The precise simulation of light scattering is a cornerstone in the development of advanced nanophotonic and biomedical devices, where functionality often hinges on the interaction between electromagnetic waves and complex, sub wavelength structures. Historically, analytical techniques like the Mie series have served as the primary tool for such analysis \citep{Wriedt2012}. However, as detailed by \citet{Martin1993}, the applicability of these formulae is rigorously constrained to perfect spheres or infinite cylinders. This introduces significant errors not only when applied to the anisotropic, swelling microstructures found in hydrogel biosensors but also when modeling irregular biological organelles or non-canonical nanostructures. To address these geometric limitations, numerical partial differential equation solvers are utilized, though each approach presents distinct trade-offs. The Finite Difference Time Domain (FDTD) method is often cited for its versatility; \citet{Drezek2000} demonstrate its utility in modeling light scattering from heterogeneous biological cells, where it offers greater geometric flexibility than Mie theory. Furthermore, \citet{teixeira2023finite} emphasize FDTD's robustness in handling broad spatial scales and multiphysics scenarios. However, both FDTD and the Finite Difference Frequency Domain (FDFD) method rely on orthogonal Cartesian grids. As \citet{Shin2013} highlights, this discretization inherently suffers from ``staircasing'' errors when modeling curved interfaces, potentially leading to inaccurate near field predictions for plasmonic nanostructures or deforming sensor surfaces. Similarly, while Boundary Element Methods (BEM) provide rigorous solutions for radiation problems \citep{Cools2011}, they can become computationally expensive when applied to penetrable objects with highly inhomogeneous material properties, such as biological tissues or gradient index lenses \citep{Buffa2003}. In this context, the Finite Element Method (FEM) offers a distinct advantage for modeling these complex, curved boundaries through the use of unstructured tetrahedral meshes, which naturally conform to both the deforming geometry of biosensors and the intricate morphology of cellular structures.

\vspace{12pt}

The choice between time domain and frequency domain formulations depends heavily on the specific optical setup and the nature of the excitation source. For applications requiring broadband spectral analysis such as analyzing the scattering cross section of a nanoparticle over the entire visible spectrum time domain methods like finite element time domain method are advantageous as they resolve the frequency response in a single simulation run \citep{teixeira2023finite}. However, for specific classes of devices interrogated by monochromatic light, the frequency domain numerical approach becomes computationally more efficient. This is particularly relevant for the stimuli responsive biosensors discussed in this study, as well as for designing specialized photonic crystals and metamaterials operating at fixed laser wavelengths. Using a time domain numerical solver for such steady state, monochromatic problems require simulating wave propagation over thousands of time steps until convergence is achieved, a process \citet{Drezek2000} note can be resource intensive. In contrast, a frequency domain numerical solver provides a direct, one-shot solution to the time harmonic Maxwell's equations, matching the continuous wave nature of the excitation. To ensure accuracy in this regime, however, the solver must be carefully designed to avoid artificial errors. Standard nodal based finite element methods often can violate the continuity of tangential electric fields across dielectric boundaries, leading to spurious solutions. \citet{Schwarzbach2009} demonstrates that this flaw frequently leads to spurious solutions that look like real signals but are numerical noise. To prevent this, it is essential to use high order Nédélec (edge) elements, which enforce the correct physical continuity of the electric field, thereby distinguishing true electromagnetic effects \citep{Koba2012} from simulation errors.

\vspace{12pt}

Finally, while the theoretical advantages of the finite element method are clear, a significant practical barrier remains due to the lack of a dedicated open source solver specifically architected for electromagnetic scattering. As highlighted by \citet{Owusu-Agyemang2025}, the scalable solution of Maxwell's equations in complex 3D geometries faces persistent challenges ranging from prohibitive computational costs to ill conditioned system matrices that require advanced interventions like parallel preconditioners to resolve. Consequently, the main problem aimed to be solved by this project is the prohibitive development curve faced by researchers attempting to utilize raw finite element frameworks, where they must manually formulate the variational weak forms and implement complex high order Nédélec (edge) elements to suppress spurious modes. This manual implementation is historically prone to error because, as established by \citet{Nedelec1980}, these vector elements require intricate covariant Piola transformations to map the tangential continuity requirements from a reference element to the distorted physical mesh. To address this bottleneck, this project leverages the FEniCS computing platform, which \citet{Otto2012} demonstrate offers a transformative advantage through its Unified Form Language (UFL). This high level abstraction allows the governing time harmonic Maxwell's equations to be expressed directly as symbolic variational forms, enabling the software to automatically handle the rigorous assembly of edge elements and their associated coordinate mappings. By building upon this advanced architecture, the project bypasses the manual coding errors identified by \citet{Boyse1992} and \citet{Jiang1996} as the primary source of numerical instability. This capability effectively removes the barrier of numerical implementation, empowering researchers to shift their focus from the tedious debugging of solver routines to the creative design and physical analysis required for next generation nanophotonic and biomedical devices.

\section{Research Objectives}

The overarching goal of this research is to develop and validate a robust finite element solver for light scattering simulations. The specific objectives are:

\begin{itemize}
    \item \textbf{Mathematical Formulation and Approximation:} To formulate the variational form of the time harmonic equation.
    
    \item \textbf{Solver Development and Stabilization:} To implement a high-accuracy finite element solver utilising unstructured meshes and vector edge elements (Nédélec elements) \citep{Nedelec1980} to ensure physical correctness and geometric flexibility.
    
    \item \textbf{Domain Truncation and Validation:} To validate robust domain truncation techniques, specifically absorbing boundary conditions (ABC), to simulate an unbounded free-space environment within a finite computational domain.
\end{itemize}

\section{Scope of Work}

This project focuses on the development of a finite element solver for Maxwell's equations in linear, isotropic, and non-magnetic media. The analysis is strictly confined to the frequency domain, assuming a time-harmonic dependence ($e^{j\omega t}$) for all fields. This approach is chosen to directly compute the steady-state scattering response without the need for time-stepping algorithms. To rigorously handle the open boundary problem, the solver utilizes the Scattered Field Formulation, which decomposes the total field into a known analytic incident wave and a numerically computed scattered wave. This specific formulation establishes the theoretical boundary for the project, necessitating a rigorous examination of the variational weak forms required to maintain the ellipticity of the system and ensure the uniqueness of the solution as discussed by \citet{Jiang1996} and \citet{Boyse1992}.

\vspace{12pt}

The geometric scope includes the simulation of scattering from dielectric and metallic cylinders with arbitrary 2D cross-sections. However, the discretization methodology is strictly defined by the requirements of high-frequency stability. Guided by the error analysis of \citet{Melenk2024}, this project prioritizes High-Order (arbitrary degree) Nédélec edge elements on unstructured triangular meshes. This choice defines the scope of the numerical investigation, shifting the focus from traditional h-refinement strategies to p-refinement techniques necessary to suppress the pollution effect inherent to indefinite Helmholtz problems. Furthermore, the reliance on $H(\text{curl})$-conforming elements to enforce tangential continuity implies a strict exclusion of nodal-based methods, which \citet{Nedelec1980} identifies as the primary source of spurious modes in electromagnetic analysis.

\vspace{12pt}

Computationally, the project involves assembling global stiffness and mass matrices to form a linear system of equations ($Ax = b$). Because the frequency-domain Helmholtz equation yields an indefinite and complex-symmetric system, the scope of the linear algebra implementation is limited to sparse direct solvers (specifically LU decomposition). As highlighted by \citet{Owusu-Agyemang2025}, this system characteristic explicitly excludes standard iterative methods (like Conjugate Gradient) which are generally unstable for wave propagation problems. Consequently, the computational framework is built upon the FEniCS platform to leverage its automated form compilation and interface with robust direct factorization libraries, as validated by \citet{Otto2012}.

\vspace{12pt}

Validation of the solver is strictly limited to the optical and microwave frequency regimes, utilizing the Analytical Mie Series for a dielectric cylinder as the sole benchmark for verification. As established by \citet{Martin1993}, this exact closed-form solution serves as the absolute reference point for verifying solver fidelity. Consequently, the validation scope focuses exclusively on the mathematical derivation of cylindrical harmonics to establish precise error metrics for scattering cross-sections and near-field intensity. Alternative semi-analytical approaches, such as the T-Matrix method, are explicitly excluded from this study to ensure that all numerical results are benchmarked against an exact, analytical standard. Nonlinear optical effects, time-varying material properties, and transient time-domain phenomena are similarly excluded.

\end{document}
